\documentclass[11pt]{article}
\usepackage[margin=1in]{geometry}
\usepackage{amsmath,amssymb,amsthm,mathtools}
\usepackage{graphicx}
\usepackage{microtype}
\usepackage[hidelinks]{hyperref}
\usepackage{enumitem}

\newtheorem{theorem}{Theorem}
\newtheorem{lemma}{Lemma}
\newtheorem{remark}{Remark}
\newcommand{\R}{\mathbb R}
\newcommand{\C}{\mathbb C}
\newcommand{\Z}{\mathbb Z}
\newcommand{\T}{\mathbb T}
\newcommand{\cH}{\mathcal H}
\newcommand{\cK}{\mathcal K}
\newcommand{\Zak}{\mathcal Z_q}
\newcommand{\one}{\mathbf 1}
\newcommand{\Span}{\overline{\operatorname{span}}}

\title{An Adjoint Gabor System Which Is Not an R-Dual of Type I}
\author{Candidate counterexample for arXiv:1509.06263\\
Run \texttt{fa\_banach\_001}, lane 6}
\date{11 August 2026}

\begin{document}
\maketitle

\begin{center}
\fbox{\parbox{0.91\textwidth}{\textbf{Status.}
Candidate counterexample, likely valid, pending specialist human review.
For every integer oversampling factor $q\ge2$, the construction gives a Gabor
frame whose normalized adjoint system is not an R-dual of type I.}}
\end{center}

\section{Source question}

Stoeva and Christensen ask whether the normalized adjoint system of every
Gabor frame is an R-dual of type I \cite[p.~6, after Theorem~1.4]{SC}.
With
\[
 (T_yf)(x)=f(x-y),\qquad (E_\xi f)(x)=e^{2\pi i\xi x}f(x),
\]
their question is whether, whenever
\(\{E_{mb}T_{na}g\}_{m,n\in\Z}\) is a frame for $L^2(\R)$, the Riesz sequence
\[
 \left\{\frac1{\sqrt{ab}}E_{m/a}T_{n/b}g\right\}_{m,n\in\Z}
\]
must be an R-dual of type I.  The source records positive answers for tight
Gabor frames and Gabor Riesz bases, but leaves the general case open.

\begin{figure}[ht]
  \centering
  \includegraphics[width=\textwidth]{figures/open_problem_crop.png}
  \caption{The open question in arXiv:1509.06263, PDF page 6.}
  \label{fig:question}
\end{figure}

The source also gives the criterion needed below.  If $S$ is the frame
operator of a frame and $S_\Omega$ is the frame operator of a proposed Riesz
sequence on its closed span, then type-I R-duality requires an antiunitary
map $\Lambda$ satisfying
\begin{equation}\label{eq:antiunitary-necessary}
 S_\Omega=\Lambda S\Lambda^{-1}.
\end{equation}
Indeed, this is the necessary direction of Lemma~1.3(iii) in \cite{SC}.

\section{Proof intuition}

At integer oversampling, there is a Parseval Gabor window whose normalized
adjoint system occupies only one of $q$ Zak components.  The adjoint shifts
act as scalar coordinate multipliers on all $q$ components.  We deform the
Parseval window by a positive invertible multiplier $\phi$ which has simple
spectrum on the two-dimensional Zak base.

The deformation makes the primal frame operator the $q$-fold amplification
$M_\phi\otimes I_q$, while the adjoint frame operator is the one-copy
operator $M_\phi$.  Choosing $\phi$ essentially one-to-one makes the second
operator cyclic.  The first operator has no cyclic vector when $q\ge2$.
Since antiunitary equivalence preserves cyclicity, the necessary condition
\eqref{eq:antiunitary-necessary} fails.

\section{The Zak model}

Fix an integer $q\ge2$ and put
\[
 a=1,\qquad b=\frac1q,\qquad
 h=q^{-1/2}\one_{[0,1)}.
\]

\begin{lemma}\label{lem:base-systems}
The Gabor system
\[
 \{E_{m/q}T_nh\}_{m,n\in\Z}
\]
is a Parseval frame for $L^2(\R)$.  Its normalized adjoint system is
\begin{equation}\label{eq:adjoint-h}
 \{\sqrt q\,E_mT_{qn}h\}_{m,n\in\Z}
 =\{E_m\one_{[qn,qn+1)}\}_{m,n\in\Z},
\end{equation}
which is an orthonormal sequence.
\end{lemma}

\begin{proof}
For each $n$, the atoms are supported on $[n,n+1)$.  Write $m=qk+j$ with
$0\le j<q$.  For fixed $j$, the functions
\[
 e^{2\pi i(k+j/q)x}\one_{[n,n+1)}(x),\qquad k\in\Z,
\]
form an orthonormal basis of $L^2([n,n+1))$.  Multiplication by
$q^{-1/2}$ and summation over the $q$ residue classes gives Parseval's
identity.  Formula \eqref{eq:adjoint-h} follows from $ab=1/q$; within each
unit interval it is the Fourier orthonormal basis, and the intervals indexed
by $n$ are disjoint.
\end{proof}

Let $X=[0,1)^2$ with coordinates $(t,\theta)$, and let
$e_0,\ldots,e_{q-1}$ be the standard basis of $\C^q$.  Define the $q$-step
Zak transform by
\begin{equation}\label{eq:zak}
 (\Zak f)_r(t,\theta)
 =\sum_{k\in\Z}f(t+r+qk)e^{-2\pi ik\theta},
 \qquad 0\le r<q.
\end{equation}
Partitioning $\R$ into the sets $t+r+qk$ and applying Fourier-series
Parseval in $k$ shows that
\[
 \Zak:L^2(\R)\longrightarrow L^2(X;\C^q)
\]
is unitary.  Direct substitution gives
\begin{equation}\label{eq:zak-shifts}
 \Zak E_1\Zak^{-1}=M_{e^{2\pi it}}\otimes I_q,
 \qquad
 \Zak T_q\Zak^{-1}=M_{e^{-2\pi i\theta}}\otimes I_q.
\end{equation}
Moreover,
\begin{equation}\label{eq:zak-h}
 \Zak h=q^{-1/2}e_0.
\end{equation}
Under $\Zak$, the atoms in \eqref{eq:adjoint-h} become
\[
 e^{2\pi imt}e^{-2\pi in\theta}e_0,
 \qquad m,n\in\Z.
\]
Consequently their closed span is
\begin{equation}\label{eq:K}
 \cK=L^2(X)e_0\subset L^2(X;\C^q).
\end{equation}

\section{Construction of the window}

We need a bounded scalar multiplier with simple spectrum.  Give each
$t,\theta\in[0,1)$ its binary expansion which is not eventually all ones,
\[
 t=.t_1t_2t_3\ldots,\qquad
 \theta=.s_1s_2s_3\ldots,
\]
and interleave the digits:
\begin{equation}\label{eq:phi}
 \psi(t,\theta)
 =\sum_{j\ge1}\left(t_j2^{-(2j-1)}+s_j2^{-2j}\right),
 \qquad \phi=1+\psi.
\end{equation}
Then $\phi:X\to[1,2)$ is a Borel injection.  Its inverse on its range is
Borel: the odd and even binary digits recover $t$ and $\theta$.

On $L^2(\R)$ define the positive invertible operator
\begin{equation}\label{eq:A}
 A=\Zak^{-1}(M_\phi\otimes I_q)\Zak
\end{equation}
and the window
\begin{equation}\label{eq:g}
 g=A^{1/2}h.
\end{equation}
Equivalently, this is the explicit Zak-coordinate prescription
\begin{equation}\label{eq:zg}
 \Zak g=q^{-1/2}\sqrt\phi\,e_0.
\end{equation}

The two unitaries in \eqref{eq:zak-shifts} generate all scalar multiplication
operators on $L^2(X;\C^q)$ by their joint bounded Borel functional calculus.
Thus $A$ belongs to the von Neumann algebra generated by $E_1$ and $T_q$.
These two generators commute with one another.  They also commute with every
primal-lattice shift $E_{m/q}T_n$, since
\[
 E_1T_n=T_nE_1,
 \qquad
 T_qE_{m/q}=E_{m/q}T_q
 \quad(m,n\in\Z).
\]
It follows that $A$ commutes with the entire primal lattice.  Because the
adjoint shifts are the products $E_mT_{qn}$, $A$ commutes with them as well.

\begin{theorem}\label{thm:counterexample}
For every integer $q\ge2$, the window $g$ in \eqref{eq:g} has the following
properties:
\begin{enumerate}[label=\textup{(\roman*)}]
\item $\{E_{m/q}T_ng\}_{m,n\in\Z}$ is a Gabor frame for $L^2(\R)$;
\item its normalized adjoint system
\[
 \{\sqrt q\,E_mT_{qn}g\}_{m,n\in\Z}
\]
is not an R-dual of type I.
\end{enumerate}
Thus the universal type-I R-duality question in \cite{SC} has a negative
answer.
\end{theorem}

\begin{proof}
Let $S_g$ be the frame operator of the primal system.  Since $A^{1/2}$
commutes with every primal shift and the system generated by $h$ is Parseval,
Lemma~\ref{lem:base-systems} gives
\begin{equation}\label{eq:Sg}
 S_g=A^{1/2}I A^{1/2}=A.
\end{equation}
In particular the system generated by $g$ is a frame, because
$I\le A\le2I$.

Write
\[
 \omega^h_{m,n}=\sqrt q\,E_mT_{qn}h,
 \qquad
 \omega^g_{m,n}=\sqrt q\,E_mT_{qn}g.
\]
Commutation with the adjoint shifts gives
\begin{equation}\label{eq:adj-deform}
 \omega^g_{m,n}=A^{1/2}\omega^h_{m,n}.
\end{equation}
The $\omega^h_{m,n}$ form an orthonormal basis of $\cK$.  Under $\Zak$,
$A^{1/2}$ is scalar multiplication by $\sqrt\phi$ and hence maps $\cK$
bijectively onto itself.  Therefore the closed span of the
$\omega^g_{m,n}$ is again $\cK$, and their frame operator on $\cK$ is
\begin{equation}\label{eq:Somega}
 S_\Omega=A^{1/2}P_\cK A^{1/2}\big|_\cK
 =A\big|_\cK.
\end{equation}
Using \eqref{eq:A}, \eqref{eq:K}, \eqref{eq:Sg}, and
\eqref{eq:Somega}, the two operators whose antiunitary equivalence is
required by \eqref{eq:antiunitary-necessary} are thus
\begin{equation}\label{eq:two-operators}
 S_g\simeq M_\phi\otimes I_q
 \quad\text{on }L^2(X;\C^q),
 \qquad
 S_\Omega\simeq M_\phi
 \quad\text{on }L^2(X).
\end{equation}

We distinguish them by cyclicity.  A vector $u$ is cyclic for a self-adjoint
operator $B$ when the closed span of
$\{F(B)u:F\text{ bounded Borel}\}$ is the whole Hilbert space.  Since
$\phi$ is a Borel injection with Borel inverse on its range, every bounded
measurable function on $X$ is a bounded Borel function of $\phi$.  Hence the
constant function $1$ is cyclic for $M_\phi$ on $L^2(X)$.

In contrast, $M_\phi\otimes I_q$ has no cyclic vector when $q\ge2$.
Indeed, given $v\in L^2(X;\C^q)$, if $v=0$ there is nothing to prove; otherwise
choose a nonzero measurable vector field $w\in L^2(X;\C^q)$ with
$w(x)\perp v(x)$ almost everywhere.  For example, using the first two
coordinates one may rotate a nonzero coordinate pair pointwise, with a
measurable partition according to the first nonzero coordinate.  Then
\[
 \langle F(\phi)v,w\rangle=0
\]
for every bounded Borel $F$, so $v$ is not cyclic.

Unitary or antiunitary equivalence of self-adjoint operators carries spectral
projections to spectral projections and therefore preserves the existence of
a cyclic vector.  The two operators in \eqref{eq:two-operators} cannot be
antiunitarily equivalent.  The necessary type-I condition
\eqref{eq:antiunitary-necessary} fails, proving (ii).
\end{proof}

\begin{remark}
The frame is neither tight nor a Riesz basis.  It is not tight because
$\phi$ is nonconstant.  The original Parseval frame is the union of $q$
orthonormal bases on each translation interval and is redundant; applying
the invertible operator $A^{1/2}$ preserves this redundancy.  Thus the
example lies genuinely outside the two positive classes recorded in the
source.
\end{remark}

\section{Verification, scope, and novelty}

The proof is exact and has no computational dependency.  The normalization
of both Gabor systems, the two Zak-shift formulas, all commutation relations,
the two frame-operator identities, and the cyclicity obstruction were checked
separately in the companion \texttt{verification.md}.  The construction uses
only a bounded positive invertible operator and therefore produces a valid
$L^2$ window.

The theorem answers the printed universal question in the negative.  Its
window is defined by a merely Borel Zak multiplier.  It does not decide
whether a type-I statement might remain true under additional regularity or
time--frequency localization hypotheses on $g$.

Before promotion, the run's registry, solution, attempt, and proof-gap indexes
were searched using arXiv:1509.06263, the title, the authors, and combinations
of \emph{R-dual of type I}, \emph{adjoint Gabor system}, \emph{counterexample},
\emph{Zak}, and spectral-multiplicity terms.  Focused arXiv/web searches
through 11 August 2026 found the source, its predecessor arXiv:1404.3183, and
the 2024 paper \cite{BDA}, but no resolution or counterexample.  The latter
still describes the question as a famous problem and develops weak R-duality
as a possible route.

Novelty confidence is moderate to high rather than conclusive: the bounded
search found no prior answer, but a specialist search of the broader Gabor
and noncommutative-torus literature is recommended.  Human verification
should focus on the joint Borel functional-calculus step and the application
of Lemma~1.3(iii) of \cite{SC}.

\begin{thebibliography}{9}
\bibitem{SC}
D.~T. Stoeva and O. Christensen,
\emph{On various R-duals and the duality principle},
Operators and Matrices 10 (2016), no.~2, 371--385;
arXiv:1509.06263.
\url{https://arxiv.org/abs/1509.06263}.

\bibitem{BDA}
H. Bansal, P. Devaraj, and S. Arati,
\emph{Characterizations of weak R-duality and its application to Gabor frames},
arXiv:2408.14952 (2024).
\url{https://arxiv.org/abs/2408.14952}.
\end{thebibliography}

\end{document}
